\documentclass[11pt,a4paper]{article}
\usepackage[top=1.7cm,bottom=1.5cm,left=1.5cm,right=1.5cm,headsep=8pt,footskip=16pt]{geometry}
\usepackage{enumitem}
\usepackage{amssymb}
\usepackage{array}
\usepackage[table]{xcolor}
\usepackage{fancyhdr}
\usepackage{lastpage}
\usepackage{multicol}
\usepackage{polyglossia}
\setdefaultlanguage[calendar=gregorian, numerals=maghrib]{arabic}
\setotherlanguage{english}
\IfFontExistsTF{Amiri}{%
  \newfontfamily\arabicfont[Script=Arabic, Scale=1.08]{Amiri}
}{%
  \newfontfamily\arabicfont[Script=Arabic, Scale=1.08, Path=../fonts/, Extension=.ttf,
    UprightFont=*-Regular, BoldFont=*-Bold, ItalicFont=*-Italic,
    BoldItalicFont=*-BoldItalic]{Amiri}
}
\IfFontExistsTF{AlBattar}{%
  \newfontfamily\khatt[Script=Arabic, Scale=1.15]{AlBattar}
}{%
  \newfontfamily\khatt[Script=Arabic, Scale=1.15, Path=../fonts/,
    BoldFont=ae_AlBattar.ttf]{ae_AlBattar.ttf}
}
\renewcommand{\headrulewidth}{0.8pt}
\renewcommand{\footrulewidth}{0.8pt}
\fancyhf{}
\rhead{\khatt المدرسة العليا للأساتذة ببوزريعة}
\lhead{\khatt ملحقة بومرداس}
\lfoot{\khatt الأستاذ: د. بن عبد الرزاق يوسف}
\cfoot{\thepage/\pageref{LastPage}}
\pagestyle{fancy}
\parindent 0pt
\setlength{\columnsep}{0.6cm}
\setlength{\columnseprule}{0.4pt}
\setlist[enumerate,1]{label=\textbf{\arabic*.}, leftmargin=*, itemsep=5pt, topsep=2pt}
\begin{document}

\linespread{1.08}\selectfont
\begin{center}
{\LARGE\khatt\textbf{الاختبار الخامس في الذكاء الاصطناعي}}\\[4pt]
{\normalsize وضعيات إشكالية ــ أسئلة متعددة الاختيارات ــ الفصول من 1 إلى 6}
\end{center}
\vspace{-2pt}\rule{\textwidth}{0.6pt}\\[2pt]
\textbf{التعليمة:} يتضمن الاختبار خمس وضعيات إشكالية. اقرأ كل وضعية بعناية، ثم اختر التحليل أو التصرف الأنسب
من بين الاختيارات الخمسة؛ لكل وضعية إجابة صحيحة واحدة فقط.
اكتب \textbf{حرف} الإجابة الصحيحة في جدول الإجابات في آخر الاختبار.
لكل وضعية 4 نقاط (المجموع 20 نقطة).
\vspace{4pt}

\noindent
{\large\khatt\textbf{الوضعية 1}}\hfill{\small(الفصل 1)}\par\nopagebreak\vspace{3pt}
\setlength{\fboxsep}{7pt}\colorbox{gray!10}{\parbox{\dimexpr\linewidth-14pt}{اقترح مدير متوسطة اعتماد أداة ذكاء اصطناعي لتصحيح أوراق الاختبار النهائي في اللغة العربية ورصد العلامات مباشرة في الكشوف، بحجة أنها أسرع وأكثر حيادًا من الأساتذة. ويتضمن الاختبار أسئلة نحو مغلقة، ووضعية إدماجية تتطلب إنتاجًا كتابيًا. وعدد الأوراق يتجاوز ستمئة ورقة، والأساتذة مثقلون بالعمل. ما التقدير الأصوب لهذا الاقتراح؟}}\par\nopagebreak
\begin{itemize}[leftmargin=1.0cm, labelsep=6pt, itemsep=3pt, topsep=4pt, parsep=0pt]
  \item[\textbf{أ)}] الاقتراح مقبول كاملًا، لأن الأداة لا تتأثر بالتعب ولا بمعرفة التلاميذ، فهي محايدة بالضرورة، وحيادها يجعل علاماتها أعدل من علامات الأساتذة في كل أنواع الأسئلة، فلا حاجة إلى مراجعة بشرية.
  \item[\textbf{ب)}] يمكن اعتماد الأداة بشرط أن يراجع الأستاذ عينة عشوائية من عشر أوراق، فإن وجدها سليمة رُصدت علامات بقية الأوراق آليًا، لأن صحة العينة تثبت صحة التصحيح في كل الأوراق مهما اختلفت إجابات التلاميذ، وهو إجراء معمول به في بعض المؤسسات.
  \item[\textbf{ج)}] يمكن أن تقترح الأداة تصحيحًا أوليًا، خاصة في الأسئلة المغلقة، وفق شبكة تقويم واضحة، لكن العلامة النهائية قرار تربوي يبقى للأستاذ بعد المراجعة، لأن تقدير الإنتاج الكتابي يتطلب حكمًا قد تخطئ فيه الأداة، وحيادها غير مضمون.
  \item[\textbf{د)}] الأداة أنسب للوضعية الإدماجية من الأسئلة المغلقة، لأن الإنتاج الكتابي يحتاج إلى تقدير لغوي تتفوق فيه الأدوات التوليدية، أما الأسئلة المغلقة فيصححها الأستاذ، ثم تُرصد العلامات دون مراجعة مشتركة.
  \item[\textbf{هـ)}] الاقتراح مرفوض كليًا، لأن الأدوات التوليدية لا تستطيع تقديم أي ملاحظة مفيدة على إجابات التلاميذ، ولا يمكن الاستفادة منها في أي مرحلة من التصحيح، حتى في اقتراح تصحيح أولي يراجعه الأستاذ، ولذلك لا معنى لإدخالها إلى المؤسسة.
\end{itemize}
\par\vspace{6pt}
\noindent
{\large\khatt\textbf{الوضعية 2}}\hfill{\small(الفصل 2)}\par\nopagebreak\vspace{3pt}
\setlength{\fboxsep}{7pt}\colorbox{gray!10}{\parbox{\dimexpr\linewidth-14pt}{لدى فريق بحث خياران لتدريب نموذج يقترح تصويبات نحوية لنصوص طلبة الجامعة: مليون جملة من المنتديات لم تُدقَّق وتكثر فيها الأخطاء، أو مئة ألف جملة دقّقها لغويون لكنها كلها من نصوص أدبية قديمة. وميزانية الفريق تسمح بتدقيق كمية محدودة إضافية. وعليه أن يقرر قبل نهاية الشهر. أي تقدير هو الأدق؟}}\par\nopagebreak
\begin{itemize}[leftmargin=1.0cm, labelsep=6pt, itemsep=3pt, topsep=4pt, parsep=0pt]
  \item[\textbf{أ)}] الخيار الأول أنسب، لأن الكمية الكبيرة تجعل الأخطاء العشوائية تُلغي بعضها بعضًا، فيتعلم النموذج الصواب من الأنماط الغالبة، وكثرة البيانات تعوّض دائمًا نقص التدقيق مهما كانت نسبة الأخطاء فيها.
  \item[\textbf{ب)}] لكل خيار عيب: الأول يعلّم النموذج أخطاء شائعة على أنها صواب، والثاني دقيق لكنه بعيد عن لغة الطلبة المعاصرة؛ والأنسب بناء مدونة مدقّقة تمثل نصوص الطلبة ولو كانت أصغر، لأن الجودة والتمثيل معًا يحددان أداء النموذج.
  \item[\textbf{ج)}] لكل خيار عيب، والحل تدريب نموذجين منفصلين على الخيارين ثم اختيار النموذج الأعلى دقة على بيانات تدريبه، لأن الدقة على بيانات التدريب هي أفضل مؤشر على أداء النموذج مع نصوص الطلبة الجديدة، ثم نشره للطلبة.
  \item[\textbf{د)}] الأنسب الجمع بين الخيارين بنسبة تسعة إلى واحد لصالح جمل المنتديات، لأن الجمل المدقّقة القليلة تكفي لتصحيح ما يتعلمه النموذج من الأخطاء الكثيرة، دون حاجة إلى تدقيق جمل المنتديات أو بناء مدونة جديدة، وهو حل اقتصادي يناسب ميزانية الفريق المحدودة.
  \item[\textbf{هـ)}] الخيار الثاني أنسب دون تحفظ، لأن التدقيق اللغوي هو المعيار الوحيد لجودة البيانات، ولا أثر لاختلاف نوع النصوص بين الأدب القديم ولغة الطلبة، فالقواعد النحوية واحدة في كل العصور والسياقات.
\end{itemize}
\par\vspace{6pt}
\noindent
{\large\khatt\textbf{الوضعية 3}}\hfill{\small(الفصل 3)}\par\nopagebreak\vspace{3pt}
\setlength{\fboxsep}{7pt}\colorbox{gray!10}{\parbox{\dimexpr\linewidth-14pt}{طلب أستاذ: «اذكر أوزان الفعل الثلاثي المزيد مع معانيها وأمثلة»، فجاءت الإجابة فقرات طويلة اختلطت فيها أوزان المزيد بحرف وبحرفين وبثلاثة، وأمثلتها نادرة الاستعمال. وهو يريد جدولًا موجزًا لتلاميذ المتوسط يُنسخ على السبورة. ويحتاج الجدول للحصة القادمة، ولا يريد أن يعيد صياغته بنفسه. أي طلب معدَّل يعالج المشكلات كلها؟}}\par\nopagebreak
\begin{itemize}[leftmargin=1.0cm, labelsep=6pt, itemsep=3pt, topsep=4pt, parsep=0pt]
  \item[\textbf{أ)}] «قدّم في جدول من أربعة أعمدة (الوزن، عدد الحروف المزيدة، المعنى الغالب، مثال شائع في جملة قصيرة) أوزان الثلاثي المزيد مرتبة حسب عدد الحروف المزيدة، بأمثلة مألوفة لتلاميذ المتوسط، دون فقرات شرح.»
  \item[\textbf{ب)}] «قدّم في جدول أوزان الثلاثي المزيد بأمثلة مألوفة لتلاميذ المتوسط، دون فقرات شرح»، لأن تحديد الأعمدة والترتيب يقيّد الأداة ويجعل الجدول أقل دقة من جدول تختار الأداة أعمدته وترتيبه بنفسها.
  \item[\textbf{ج)}] «قدّم أوزان الثلاثي المزيد في جدول مرتب ومنظم وجميل يناسب السبورة، دون فقرات شرح»، لأن ذكر الجدول وحده يكفي لتنظيم الأوزان حسب عدد الحروف واختيار أمثلة مألوفة تلقائيًا.
  \item[\textbf{د)}] «قدّم في جدول من أربعة أعمدة (الوزن، عدد الحروف المزيدة، المعنى الغالب، مثال) أوزان الثلاثي المزيد مرتبة حسب عدد الحروف المزيدة، ثم أضف بعد الجدول شرحًا مفصلًا لكل وزن بأمثلة من الشعر القديم.»
  \item[\textbf{هـ)}] «أعد كتابة الإجابة السابقة في جدول» فقط، لأن الأداة تحتفظ بالمحتوى السابق، فيكفي تغيير الشكل لتصبح الأوزان مرتبة والأمثلة مألوفة، دون ذكر المستوى أو الأعمدة أو طريقة الترتيب.
\end{itemize}
\par\vspace{6pt}
\noindent
{\large\khatt\textbf{الوضعية 4}}\hfill{\small(الفصل 5)}\par\nopagebreak\vspace{3pt}
\setlength{\fboxsep}{7pt}\colorbox{gray!10}{\parbox{\dimexpr\linewidth-14pt}{لدى أستاذة تقرير من ثلاثين صفحة عن المطالعة، يتضمن نتائج استبيان، وتوصيات موجهة للأسرة، وأخرى موجهة للمؤسسة. تريد ملخصًا لتلاميذ السنة الأولى متوسط، وآخر تعرضه على أوليائهم في اجتماع مدته عشر دقائق. ووقتها لا يسمح بقراءة التقرير مرات عديدة. ما الطريقة الأنسب لاستعمال أداة الذكاء الاصطناعي في هذه المهمة؟}}\par\nopagebreak
\begin{itemize}[leftmargin=1.0cm, labelsep=6pt, itemsep=3pt, topsep=4pt, parsep=0pt]
  \item[\textbf{أ)}] تطلب ملخصًا واحدًا شاملًا ثم تطلب تبسيطه للتلاميذ، وتقدّم الملخص الشامل للأولياء، لأنه يصلح لكل البالغين، وتكتفي بمراجعة الملخص المبسّط لأن الملخص الشامل مأخوذ مباشرة من التقرير.
  \item[\textbf{ب)}] تطلب ملخصين منفصلين بتحديد الجمهور والطول ومستوى اللغة، وتعتمدهما دون مقابلتهما بالتقرير، لأن تحديد هذه العناصر يضمن أمانة الملخص، فالأداة لا تحرّف النتائج إلا عندما يكون الطلب غامضًا، وهذا يوفر وقتها.
  \item[\textbf{ج)}] تقسم التقرير إلى ثلاثة أجزاء وتلخص كل جزء وحده، ثم تجمع الملخصات الثلاثة وتقدّمها للجمهورين كما هي، لأن تلخيص الأجزاء منفصلة يحفظ الأمانة ويغني عن تكييف الملخص مع كل جمهور.
  \item[\textbf{د)}] تطلب ملخصين منفصلين، وتطلب في ملخص الأولياء التركيز على نتائج الاستبيان وتوصيات المؤسسة لأنها الأجزاء الأكثر علمية، ثم تقابل الملخصين بالتقرير، لأن الأولياء يهتمون بالأرقام أكثر من التوصيات الموجهة إليهم، فيصبح الملخص أنفع لهم.
  \item[\textbf{هـ)}] تطلب ملخصين منفصلين تحدد لكل منهما الجمهور والطول ومستوى اللغة والأجزاء ذات الأولوية، كتوصيات الأسرة للأولياء، ثم تقابل كل ملخص بالتقرير لتتأكد أنه لم يضف معلومة ولم يحرّف نتيجة ولم يُسقط فكرة جوهرية.
\end{itemize}
\par\vspace{6pt}
\noindent
{\large\khatt\textbf{الوضعية 5}}\hfill{\small(الفصل 6)}\par\nopagebreak\vspace{3pt}
\setlength{\fboxsep}{7pt}\colorbox{gray!10}{\parbox{\dimexpr\linewidth-14pt}{تحضّر أستاذة حصة في التعبير الكتابي عن وصف مشهد طبيعي لتلاميذ السنة الثانية متوسط، وتريد استثمار أدوات متعددة الأنماط لتحفيزهم، مع الحرص على أن يبقى الإنتاج الكتابي من عمل التلاميذ، وأن يكتسبوا أساليب الوصف. وقد لاحظت في حصص سابقة أن تلاميذها يملّون من الوصف المجرد بلا سند. أي تصوّر للنشاط هو الأنسب تربويًا؟}}\par\nopagebreak
\begin{itemize}[leftmargin=1.0cm, labelsep=6pt, itemsep=3pt, topsep=4pt, parsep=0pt]
  \item[\textbf{أ)}] تعرض الصورة والمقطع الصوتي، ثم تقدّم وصفًا نموذجيًا مولَّدًا قبل الكتابة ليحاكيه التلاميذ، لأن الاطلاع على النموذج أولًا يضمن جودة النصوص، ثم تقيّم نصوصهم بمدى قربها من النموذج المولَّد.
  \item[\textbf{ب)}] تطلب من التلاميذ وصف المشهد شفهيًا لأداة صوتية تحوّل كلامهم إلى نص مكتوب منقّح، ثم تقيّم النص المنقّح على أنه إنتاجهم الكتابي، لأن الأفكار أفكارهم، والتنقيح الآلي لا يغيّر شيئًا في ملكيتهم للنص.
  \item[\textbf{ج)}] تعرض عشرات الصور المولَّدة لمشاهد مختلفة طوال الحصة، وتناقش التلاميذ شفهيًا في جمالها، لأن كثرة الصور تكفي لتنمية مهارة الوصف الكتابي، فلا حاجة إلى تخصيص وقت للكتابة أو للمقارنة.
  \item[\textbf{د)}] تعرض صورة مولَّدة لمشهد ومقطعًا صوتيًا لأصواته، فيلاحظ التلاميذ ويدوّنون مفردات حسية، ثم يكتبون وصفهم بأنفسهم، وبعدها تقارن معهم نصوصهم بوصف نموذجي راجعته مسبقًا، لتبيين أساليب الوصف لا لنسخها.
  \item[\textbf{هـ)}] تعرض الصورة والمقطع الصوتي، ويكتب التلاميذ وصفهم بأنفسهم، ثم تطلب من أداة نصية تصحيح نصوصهم وتعيدها إليهم مصححة دون مناقشة، لأن التصحيح الآلي أدق من الأستاذ ويغني عن تبيين الأساليب.
\end{itemize}
\par\vspace{6pt}
\begin{center}
{\khatt\textbf{جدول الإجابات}}\\[3pt]
\renewcommand{\arraystretch}{1.6}
\begin{tabular}{|>{\centering\arraybackslash}m{1.9cm}|*{5}{>{\centering\arraybackslash}m{1.2cm}|}}
\hline
\textbf{رقم السؤال} & \textbf{1} & \textbf{2} & \textbf{3} & \textbf{4} & \textbf{5} \\ \hline
\textbf{الإجابة} &  &  &  &  &  \\[8pt] \hline
\end{tabular}

\vspace{8pt}
\renewcommand{\arraystretch}{2.2}
\begin{tabular}{|>{\centering\arraybackslash}m{7.5cm}|>{\centering\arraybackslash}m{3.2cm}|>{\centering\arraybackslash}m{3.2cm}|}
\hline
\textbf{الاسم واللقب} & \textbf{الفوج} & \textbf{العلامة} \\ \hline
 & & \hfill /20 \\ \hline
\end{tabular}
\end{center}
\end{document}
